\documentclass[10pt,a4paper]{amsart}
\usepackage[margin=19mm]{geometry}
\usepackage{amsmath,amssymb,mathtools}
\usepackage[scr=rsfs,cal=euler]{mathalfa}
\usepackage{xfrac}
\usepackage[unicode,hidelinks,bookmarks=false]{hyperref}
\newcommand{\ie}{i.e.\ }
\newcommand{\cf}{cf.\ }
\newcommand{\resp}{respectively\ }
\newcommand{\Subsection}{\subsection}
\providecommand{\nobreakdash}{\nobreak\hbox{-}}
\setlength{\emergencystretch}{2em}
\multlinegap0pt
\usepackage{bm}
\usepackage{bbm}
\usepackage{dsfont}
\usepackage{xspace}
\usepackage{cancel}
\usepackage{mathtools}
\usepackage{amsthm}
\makeatletter
\@ifundefined{theorem}{\newtheorem{theorem}{Theorem}}{}
\makeatother
\title[Recovering Riemannian Geometry from Diffusion]{Recovering Riemannian Geometry from Diffusion}
\author{Amandip Sangha}
\date{Source: arXiv:2601.17166v1 (CC BY 4.0)}
\begin{document}
\maketitle

\noindent\emph{Source and licence.} Recovering Riemannian Geometry from Diffusion. Version arXiv:2601.17166v1 (CC BY 4.0), distributed under the Creative Commons Attribution 4.0 International licence, as recorded in the arXiv licence metadata for this exact version and shown on the abstract page https://arxiv.org/abs/2601.17166v1. Full rights notice: licenses/arxiv-2601.17166-CC-BY-4.0.txt.\par\smallskip

\noindent\textit{Editorial scope.} Counted unit: the Theorem whose source label is \texttt{thm:recover-metric-from-Gamma} in the article (source0.tex lines 676--684, source label \texttt{thm:recover-metric-from-Gamma}), delivered together with the article's own complete original proof of it (source0.tex lines 685--766, one \texttt{proof} environment of 82 lines). No mathematics is written, repaired or re-derived: statement and proof are the article's own text line for line. Required context retained uncounted: eq:cometric-def (lines 621--625). Every definition, display and label of those lines is retained unchanged. Omitted from this package and not redistributed: 1--620, 626--675, 767--1709. Nothing inside a retained excerpt is omitted or altered: every retained body is delivered exactly as the article's lines are written, so no omission receipt of any kind accompanies this package. Citations: the retained excerpts carry no citation command at all, so no bibliography entry of the article is transcribed and no reference is redistributed. Reference adaptation: none was needed and none was made; every reference inside the delivered unit resolves inside it, so no unresolved reference or citation is produced. Printed slips kept literal: no printed word, symbol, hypothesis or number of the counted statement or of its proof was altered. Layout and class: the article's own class is \texttt{amsart} with packages including amsmath, amssymb, amsthm, amsfonts, mathrsfs, hyperref, enumerate, tikz, bm, geometry, color, soul; the wrapper is the lane's pinned \texttt{amsart} editorial wrapper at 10pt on A4, which restates the theorem-like environments the retained text needs and adds the article's own macro definitions that the retained text needs (none), with extra wrapper packages (bm, bbm, dsfont, xspace, cancel, mathtools). The delivered numbering is the unit's own. Assets: no retained span contains a figure environment, a graphics-inclusion command, a PDF-page inclusion, a table or a TikZ picture; no image file is copied into this package, the package declares no asset dependency and no image file is redistributed with it. Licence: version arXiv:2601.17166v1 of this article is distributed under the Creative Commons Attribution 4.0 International licence, as recorded in the arXiv licence metadata for this exact version and shown on the abstract page https://arxiv.org/abs/2601.17166v1. A full rights notice is delivered at licenses/arxiv-2601.17166-CC-BY-4.0.txt. Characters: every retained excerpt is pure ASCII, so the delivered engine, XeLaTeX with its default Latin Modern text font, reports no missing character.
\begin{equation}\label{eq:cometric-def}
\langle \alpha,\beta\rangle_{g_x^{-1}}
\;:=\;
\Gamma(f,g)(x),
\end{equation}

\begin{theorem}[Metric recovered from $\Gamma$]\label{thm:recover-metric-from-Gamma}
Assume in addition that the bilinear form \eqref{eq:cometric-def} is positive definite at every
$x\in M$.  Then there exists a unique smooth Riemannian metric $g$ on $M$ such that for all
$f,g\in C^\infty(M)$,
\begin{equation}\label{eq:Gamma-metric}
\Gamma(f,g)=\langle df,dg\rangle_{g^{-1}}.
\end{equation}
Equivalently, in local coordinates $(x^i)$ one has $\Gamma(x^i,x^j)=g^{ij}$.
\end{theorem}

\begin{proof}
Fix $x\in M$. By Assumption~2.2, for any $f,g\in C^\infty(M)$ the value $\Gamma(f,g)(x)$ depends
only on the differentials $df(x),dg(x)\in T_x^*M$. Hence the prescription
\[
\langle \alpha,\beta\rangle_{g^{-1}_x}:=\Gamma(f,g)(x),
\qquad df(x)=\alpha,\ dg(x)=\beta,
\]
defines a \emph{well-defined} symmetric bilinear form on $T_x^*M$ (cf.\ (5)).
By hypothesis, this bilinear form is positive definite for every $x$, so $g^{-1}_x$ is an inner product
on $T_x^*M$.

Let $(U;x^1,\dots,x^n)$ be a chart and define on $U$
\[
G^{ij}:=\Gamma(x^i,x^j).
\]
Since $\Gamma$ is smooth on smooth inputs (Assumption~2.2), each $G^{ij}\in C^\infty(U)$.
Moreover, for any $f,h\in C^\infty(U)$, writing $f=F(x^1,\dots,x^n)$ and $h=H(x^1,\dots,x^n)$
locally and applying the multivariate chain rule (Lemma~2.3) in each slot yields the coordinate formula (7):
\begin{equation}\label{eq:Gamma-coords}
\Gamma(f,h)=\sum_{i,j=1}^n (\partial_i f)(\partial_j h)\,G^{ij}
\qquad\text{on }U.
\end{equation}
Evaluating \eqref{eq:Gamma-coords} at $x\in U$ and taking $f=x^i$, $h=x^j$ gives
\[
\langle dx^i,dx^j\rangle_{g^{-1}} = \Gamma(x^i,x^j)=G^{ij},
\]
so $(G^{ij}(x))$ is the matrix of $g^{-1}_x$ in the coframe $(dx^i)_x$. Positivity of $g^{-1}_x$
is equivalent to $(G^{ij}(x))$ being symmetric positive definite, hence invertible. Define its inverse
\[
(g_{ij}) := (G^{ij})^{-1}.
\]
Since inversion is smooth on the cone of positive definite matrices, $g^{ij}\in C^\infty(U)$.

Let $(V;y^1,\dots,y^n)$ be another chart with $U\cap V\neq\varnothing$. On $U\cap V$, apply
Lemma~2.3 to each coordinate function $y^a=y^a(x^1,\dots,x^n)$ and use bilinearity of $\Gamma$:
\begin{align}
G^{(y)}_{ab}
:=\Gamma(y^a,y^b)
&=\sum_{i,j=1}^n (\partial_i y^a)(\partial_j y^b)\,\Gamma(x^i,x^j)
=\sum_{i,j=1}^n (\partial_i y^a)(\partial_j y^b)\,G^{(x)}_{ij}.
\label{eq:cometric-transform}
\end{align}
Thus the local matrices $(G^{ij})$ transform on overlaps as the components of a smooth $(2,0)$-tensor,
i.e.\ a globally defined smooth co-metric $g^{-1}\in C^\infty(M;\,S^2TM)$. Since each $g^{-1}_x$ is
positive definite, the pointwise inverse defines a smooth $(0,2)$-tensor
\[
g := (g^{-1})^{-1}\in C^\infty(M;\,S^2T^*M),
\]
whose local coefficients in the chart $(x^i)$ are exactly $(g_{ij})=(G^{ij})^{-1}$, and by
\eqref{eq:cometric-transform} these coefficients patch to a global tensor. Positivity implies $g$ is a
Riemannian metric.

Fix $x\in U$ and write $df(x)=(\partial_i f)(x)\,dx^i_x$ and $dh(x)=(\partial_j h)(x)\,dx^j_x$.
Then, since $g^{-1}(dx^i,dx^j)=G^{ij}$,
\[
\langle df,dh\rangle_{g^{-1}}(x)
=\sum_{i,j=1}^n (\partial_i f)(x)(\partial_j h)(x)\,G^{ij}(x)
=\Gamma(f,h)(x)
\]
by \eqref{eq:Gamma-coords}. Hence (9) holds pointwise on $M$.

Suppose $\tilde g$ is another smooth Riemannian metric with
\[
\Gamma(f,h)=\langle df,dh\rangle_{\tilde g^{-1}}
\qquad\forall f,h\in C^\infty(M).
\]
Fix $x\in M$ and $\alpha,\beta\in T_x^*M$. Choose a chart $(U;x^1,\dots,x^n)$ about $x$ and write
$\alpha=\alpha_i\,dx^i_x$, $\beta=\beta_j\,dx^j_x$. Let $\chi\in C_c^\infty(U)$ satisfy $\chi\equiv 1$
near $x$, and define
\[
f:=\chi\sum_{i=1}^n \alpha_i x^i,\qquad h:=\chi\sum_{j=1}^n \beta_j x^j.
\]
Then $df(x)=\alpha$ and $dh(x)=\beta$. Therefore
\[
\langle \alpha,\beta\rangle_{\tilde g^{-1}_x}
=\langle df(x),dh(x)\rangle_{\tilde g^{-1}_x}
=\Gamma(f,h)(x)
=\langle df(x),dh(x)\rangle_{g^{-1}_x}
=\langle \alpha,\beta\rangle_{g^{-1}_x}.
\]
Since $\alpha,\beta$ were arbitrary, $\tilde g^{-1}_x=g^{-1}_x$ for all $x$, hence $\tilde g=g$.
\end{proof}

\end{document}
